% 05-related.tex — §5 相关工作与讨论
\section{相关工作与讨论}
\label{sec:related}

\subsection{可验证窥孔优化框架}

我们从作用的抽象层次、验证方法与进位链覆盖三个维度定位本文工作。基于 SMT 求解器校验的窥孔优化工作可追溯到 Lopes 等人提出的
Alive~\cite{alive}及其后继 Alive2~\cite{alive2}，以及 Sasnauskas 等人的
Souper~\cite{souper}——它们共同确立了“先合成候选、再用 SMT 等价性验证准入”的
基本范式。本文沿用此范式，但将校验的 IR 从通用 LLVM IR 改为 EVM JIT 的 dMIR。
值得说明的是，Alive 的 PLDI 2015 论文已明确指出 LLVM 以多返回值结构表达进位/溢出
标志对窥孔验证的表达力限制——这一观察是本文的起点，而非新发现。

Bhat 等人在 ITP 2024 提出的 LeanMLIR~\cite{leanmlir}在 Lean~4 中形式化了
“将窥孔验证框架对 IR 参数化”的思想，与本文的定位互补：LeanMLIR 关注证明框架
对 IR 的通用性，本文关注在具体 IR（dMIR）上的实例化规则集与端到端集成。

Mukherjee 与 Regehr 在 OOPSLA 2024 提出的 Hydra~\cite{hydra}可自动发现与
输入位宽无关的窥孔规则；由于 Hydra 与 Alive2 同样建立在 LLVM IR 之上，在
多字进位链这族规则上承袭第~2~节所述的表达力限制，本文通过更换 IR 基础给出
一种工程解法。Schett 与 Nagele 在 FMBC 2020 发表的工作~\cite{fmbc20_evm_peephole}
以相似的离线枚举加 SMT 流程为 EVM \emph{字节码}自动生成约 993 条规则；两者
在等价性保证上对等，差别在抽象层次：本文规则集作用在 JIT 代码生成路径上，
可同时利用 dMIR 的结构化匹配与 x86 层机械冗余清理（字节码层均不可见），
规则数量差距（993 对 83）反映作用层次与覆盖策略不同。

多字整数的 \texttt{adc}/\texttt{sbb} 进位链是 GMP、libsecp256k1 等密码学库的
规范下降形态并对应 Intel ADCX/ADOX 指令；4 段 64 位链是 EVM 256 位算术在
64 位硬件上的典型形状，并非刻意构造的边界情形。

\subsection{面向共识关键路径的验证工作}

Albert 等人于 CAV 2023 发表的工作~\cite{verified_evm_cav23}以 Coq 为证明基础，
在 EVM 字节码块级 AOT 优化规则集合上完成形式化证明。其定位是字节码块级、
离线证明、通用位向量与代数恒等；本文定位是 JIT 中间表示与机器码层、规则合成
时的 SMT 准入、包含多字进位链的 EVM 专用模式。两者规模并不对等：Albert 的
14 条规则\footnote{以 artifact 的 stack-only 分支 Coq 源文件
\texttt{optimizations.v} 为准；原文 extended version 的 Appendix A 可能列出
更多 memory/storage 层规则，两侧系统均不涉及，此处按 stack-only 口径对照。}
均作用于字节码块级的 stack 代数化简；本文 83 条与其重合 7 条，均为 verified
位向量窥孔系统几乎必然包含的代数恒等公理（加法/乘法幺元、乘零吸收、或零、
双重否定、自减、与吸收等；在 Alive2、Souper 与 LLVM InstCombine 中均可见
同构实现），不构成衍生依赖。其余 76 条中 13 条位于 x86 代码生成层，60 余条
位于 dMIR 层且覆盖 Albert 未涉及的范畴——包含 256 位字长分解出的多字进位链
规则、\texttt{adc}/\texttt{sbb} 语义的代数简化、与 JIT 代码生成路径相关的
常量传播与强度削减等。

运行时对比验证工作（如 Breidenbach 等人 USENIX Security 2018 发表的 Hydra
Framework~\cite{hydra_framework}所代表的多客户端共识检测系统、EVM 模糊测试
以及共识测试网络）在上线后检测分歧；本文定位为开发期的静态强化。

\subsection{局限与未来工作}
\label{sec:limitations}

\begin{enumerate}[leftmargin=1.8em]
  \item \textbf{可信基础。}SMT 等价性校验建立在 dMIR 位向量语义模型之上，
    我们假设该模型与真实 JIT 执行引擎的运行语义一致；这一一致性的系统化测试
    作为未来工作。
  \item \textbf{x86 代码生成层规则的校验覆盖。}13 条 x86 层规则均为机械语法
    重写，未配套独立 SMT 校验，其正确性由第~\ref{sec:correctness}~节两侧对照
    测试套件保障；x86 寄存器与标志位语义的形式化建模作为未来工作。
  \item \textbf{工作负载覆盖。}本文加速数据来自 evmone-bench 的 27 个基准，
    不能代表全部生产合约。
  \item \textbf{规则建模范围。}本文只对窥孔重写前后的算术与标志位语义做等价性
    校验，不建模气体计费、内存别名与存储槽等副作用。
  \item \textbf{死进位消除在 U256 加减法上的构造性阻塞。}U256 加减法下降出的
    链首带进位加法，其第三操作数是运行时值而非可匹配的零常量，该规则在链首按
    构造无法生效；保留它是因为它在 U256 乘法下降内部的独立带进位加法上仍有效。
  \item \textbf{评测精度与编译开销对照。}\texttt{weierstrudel/15} 与
    \texttt{signextend/one} 落在 $\pm 2\%$ 噪声带内；\texttt{weierstrudel/1}
    稳定慢 \textbf{2.71\%}，逐一关闭新增规则对照未发现稳定负贡献。单模块编译
    时间中位数 0.45\,ms、95 百分位 0.87\,ms 均来自启用侧单边测量；EVMC 接口
    无逐阶段计时钩子，基线对照与窥孔重写单步开销分离均需另行插桩，本文仅以
    启用侧作为整体编译延迟的上界。
  \item \textbf{\texttt{memory\_grow\_mstore} 族微基准回归。}4 个
    \texttt{memory\_grow\_mstore} 变体稳定慢 \textbf{7.74\%}--\textbf{8.85\%}
    （两侧 CV 均低于 2\%）。初步分析把回归定位到 x86 层比较合并规则
    \texttt{optimizeCmp} 在 memory-expand 控制流下降路径上的匹配行为——该
    路径以 \texttt{SETCC}$\to$ 零扩展$\to$\texttt{TEST}$\to$\texttt{JCC}
    链式下降，其中零扩展使用 \texttt{SUBREG\_TO\_REG} 形式，与当前规则集覆盖
    的紧凑链在模式上不完全一致；具体根因留作后续工作。同族 \texttt{mload}
    4 个变体不涉及此合并机会，加速 \textbf{+6.03\%}--\textbf{+7.26\%}。
\end{enumerate}
